% ============================================================================
% GAHIB: Graph Attention VAE with Hyperbolic Information Bottleneck
%        for Single-Cell Omics Latent Space Learning
%
% Target: Elsevier journal (elsarticle single-column)
% ============================================================================
\documentclass[preprint,12pt]{elsarticle}

% ── Packages ──────────────────────────────────────────────────────────────
\usepackage[T1]{fontenc}
\usepackage{amsmath,amssymb,amsfonts}
\usepackage{graphicx}
\usepackage{booktabs}
\usepackage{multirow}
\usepackage{xcolor}
\usepackage{hyperref}
\usepackage{cleveref}
\usepackage{siunitx}
\usepackage{microtype}
\usepackage{enumitem}
\usepackage{algorithm}
\usepackage{algorithmic}
\usepackage{lineno}
\usepackage{rotating}
\usepackage{placeins}
\usepackage[a4paper,margin=2.2cm,top=2.0cm,bottom=2.0cm]{geometry}

% Float settings tuned to reduce residual whitespace without destabilising order
\renewcommand{\topfraction}{0.95}
\renewcommand{\bottomfraction}{0.95}
\renewcommand{\textfraction}{0.05}
\renewcommand{\floatpagefraction}{0.85}
\setcounter{topnumber}{4}
\setcounter{bottomnumber}{4}
\setcounter{totalnumber}{8}

% Reduce spacing around floats
\setlength{\textfloatsep}{6pt plus 2pt minus 4pt}
\setlength{\floatsep}{3pt plus 2pt minus 2pt}
\setlength{\intextsep}{4pt plus 2pt minus 3pt}
\setlength{\dbltextfloatsep}{6pt plus 2pt minus 4pt}
\setlength{\dblfloatsep}{3pt plus 2pt minus 2pt}

% Reduce spacing around equations
\setlength{\abovedisplayskip}{4pt plus 2pt minus 2pt}
\setlength{\belowdisplayskip}{4pt plus 2pt minus 2pt}
\setlength{\abovedisplayshortskip}{2pt plus 1pt minus 1pt}
\setlength{\belowdisplayshortskip}{2pt plus 1pt minus 1pt}

% Reduce spacing around section headings (elsarticle-compatible)
\makeatletter
\renewcommand\section{\@startsection{section}{1}{\z@}%
  {-8pt \@plus -2pt \@minus -2pt}%
  {4pt \@plus 1pt \@minus 1pt}%
  {\normalfont\Large\bfseries}}
\renewcommand\subsection{\@startsection{subsection}{2}{\z@}%
  {-6pt \@plus -2pt \@minus -2pt}%
  {3pt \@plus 1pt \@minus 1pt}%
  {\normalfont\large\bfseries}}
\renewcommand\paragraph{\@startsection{paragraph}{4}{\z@}%
  {-4pt \@plus -1pt \@minus -1pt}%
  {-0.5em}%
  {\normalfont\normalsize\bfseries}}
\makeatother
\raggedbottom

% Tighten paragraph and list spacing
\setlength{\parskip}{0pt plus 1pt}
\setlength{\parsep}{0pt plus 1pt}
\setlength{\partopsep}{0pt plus 1pt}

\graphicspath{{figures/}}

% ── Math operators ────────────────────────────────────────────────────────
\DeclareMathOperator{\KL}{KL}
\DeclareMathOperator{\Var}{Var}
\DeclareMathOperator{\diag}{diag}
\newcommand{\rea}{\mathbb{R}}
\newcommand{\lore}{\mathbb{L}}
\newcommand{\poin}{\mathbb{D}}
\newcommand{\bx}{\mathbf{x}}
\newcommand{\bz}{\mathbf{z}}
\newcommand{\bmu}{\boldsymbol{\mu}}
\newcommand{\bsig}{\boldsymbol{\sigma}}

% ── Significance markers ─────────────────────────────────────────────────
\newcommand{\sigA}{$^{***}$}     % p < 0.001
\newcommand{\sigB}{$^{**}$}      % p < 0.01
\newcommand{\sigC}{$^{*}$}       % p < 0.05
\newcommand{\ns}{$^{\text{ns}}$} % not significant

\hypersetup{
  colorlinks=true,
  linkcolor=blue!60!black,
  citecolor=blue!60!black,
  urlcolor=blue!60!black,
  hypertexnames=false,
}

\journal{Journal of King Saud University -- Computer and Information Sciences}

% ============================================================================
\begin{document}

\begin{frontmatter}

\title{Graph Attention VAE with a Hyperbolic Information Bottleneck\\
for Single-Cell Latent Representation Learning}

\author[shef]{Zeyu Fu\corref{cor}}
\ead{zfu3@sheffield.ac.uk}
\cortext[cor]{Corresponding author.\\
  ORCID: \url{https://orcid.org/0009-0001-8329-0108}}
\address[shef]{Department of Computer Science, University of Sheffield,
Sheffield, United Kingdom}

% ── Abstract ──────────────────────────────────────────────────────────────
\begin{abstract}
Learning meaningful latent representations from single-cell RNA-sequencing
(scRNA-seq) data requires models that preserve both local neighbourhood
structure and hierarchical biological organisation, yet existing
variational autoencoders encode cells independently in Euclidean space
and cannot capture these properties simultaneously.
We introduce \textbf{GAHIB} (\emph{Graph Attention VAE with Hyperbolic
Information Bottleneck}), which integrates a graph-attention encoder,
a 2D information bottleneck, and a Lorentz-hyperbolic geometry loss
in one generative framework.
The GAT encoder propagates information over a cell-neighbourhood
graph; the bottleneck provides a compressed manifold coordinate; and
the hyperbolic term anchors the latent space to hierarchical structure.
We evaluate GAHIB on 53 scRNA-seq datasets through eleven studies:
seven comparative benchmarks (ablation, deep-learning, classical DR,
geometric VAE, disentanglement, encoder, graph convolution) and four
robustness studies (latent dimension, multi-seed, cost, hyperparameter).
GAHIB delivers the strongest overall clustering and
dimensionality-reduction fidelity, with statistically significant
gains over most baselines under Wilcoxon signed-rank tests.
Hyperbolic radii align with developmental hierarchy
($\rho_s{=}0.327$, up to 0.673), decoder Jacobians recover marker
genes, and GAT attention exhibits strong cell-type homophily.
GO~Biological~Process enrichment on four tissue contexts shows that
latent dimensions encode distinct, tissue-specific programmes without
supervision. Pseudotime reconstructed from Lorentz norms correlates
tightly with diffusion pseudotime. GAHIB is reproducible
(${<}$3\% std across 5 seeds), robust across hyperparameters, and
scales flatly with cell count, providing a principled, efficient,
biologically grounded framework for single-cell representation learning.
\end{abstract}

\begin{keyword}
Single-cell RNA-seq \sep graph attention network \sep
variational autoencoder \sep hyperbolic geometry \sep
information bottleneck \sep representation learning
\end{keyword}

\begin{highlights}
\item GAHIB couples graph attention, information bottleneck, and Lorentz
  hyperbolic geometry in a single single-cell VAE architecture.
\item Eleven experimental studies on 53 scRNA-seq datasets show GAHIB
  outperforms 7 deep-learning and 5 classical baselines on 20 metrics.
\item Latent hyperbolic radii track developmental hierarchy
  ($\rho_s{=}0.327$) without supervision, yielding emergent stemness
  scores.
\item Tissue-specific GO Biological Process enrichment confirms each
  latent dimension encodes biologically coherent, context-dependent
  programmes.
\item Model is reproducible ($<$3\% std across 5 seeds), robust to
  hyperparameter choices, and scales flat with cell count.
\end{highlights}

\end{frontmatter}

\linenumbers


% ============================================================================
\section{Introduction}\label{sec:intro}
% ============================================================================

\begin{figure}[!t]
\centering
\includegraphics[width=\textwidth]{overview}
% Source: paper/figures/overview.pdf
% Generator: gahib/viz/fig_overview.py
\caption{Overview of GAHIB.
\textbf{(a)}~GAHIB combines graph attention, an information bottleneck,
and hyperbolic regularisation to learn biologically structured embeddings.
\textbf{(b)}~Raw counts are converted into a cell graph, encoded into a
10-dimensional latent space, and decoded with a negative-binomial head.
\textbf{(c)}~Evaluation spans 53 scRNA-seq datasets, eleven study
tracks, and 20 metrics.}
\label{fig:overview}
\end{figure}

Single-cell RNA-sequencing (scRNA-seq) has transformed our ability to
characterise cellular heterogeneity, producing high-dimensional count
matrices that capture thousands of genes per
cell~\cite{zheng2017massively}.
Extracting compact, biologically informative latent representations from
these data underpins most downstream tasks---clustering,
trajectory inference, differential expression analysis, and cell-type
annotation---making representation learning a central objective in
computational single-cell genomics.
Widely used analysis pipelines such as
Scanpy~\cite{wolf2018scanpy} and Leiden
clustering~\cite{traag2019leiden} depend critically on the quality of
these latent embeddings.

Unlike bulk expression measurements, scRNA-seq matrices are sparse,
overdispersed, and structurally heterogeneous: discrete cell identities
coexist with continuous transitional states, branching developmental
programmes, and local neighbourhood effects induced by shared lineage or
microenvironment.
An embedding that is useful in practice therefore has to satisfy several
requirements simultaneously.
It must denoise counts, preserve local neighbour relations for clustering
and graph construction, and retain enough global organisation for
trajectory analysis and biological interpretation.
Methods that optimise only one of these goals often produce either clean
clusters with poor manifold structure or smooth manifolds with weak
cell-state separation.

Variational autoencoders (VAEs) have emerged as the \emph{de facto}
framework for single-cell representation
learning~\cite{lopez2018deep,gronbech2020scvae}.
scVI~\cite{lopez2018deep} exemplifies this paradigm, fitting a
negative-binomial likelihood with a multi-layer perceptron (MLP) encoder to
yield Euclidean embeddings suitable for downstream clustering.
These models are effective denoisers, but their default Euclidean latent
spaces still treat neighbourhood continuity and lineage branching as
secondary effects rather than explicit structural constraints.
Despite widespread adoption, standard VAEs overlook two structural priors
that are known to govern biological organisation.
First, cell types form \emph{differentiation hierarchies}
(stem $\to$ progenitor $\to$ differentiated), a branching topology that
Euclidean geometry cannot represent without bounded distortion:
hyperbolic space embeds tree-like structures with exponentially lower
distortion~\cite{nickel2017poincare,ganea2018hyperbolic}.
Second, cells form \emph{local neighbourhoods} in gene-expression space,
yet most VAE encoders treat each cell independently, discarding the graph
topology entirely.

Recent methods have begun to address these limitations, but only
individually.
Geometric VAEs~\cite{bayer2023gmvae} explore hyperbolic priors;
scGNN~\cite{wang2021scgnn} and scGCC~\cite{tian2023scgcc} build
graph-aware encoders;
scDHMap~\cite{tian2023scdhmap} maps latent codes onto the Poincar\'{e}
disk;
and siVAE~\cite{choi2023sivae} extracts gene-level interpretability via
decoder Jacobians.
The deep variational information bottleneck~\cite{alemi2017deep} provides
a variational framework for structured compression, yet has not been
coupled to hyperbolic geometry in single-cell settings.
Critically, \emph{no existing method jointly combines} graph-structured
encoding, information bottleneck compression, and hyperbolic geometry
within a unified generative framework---the gap that GAHIB addresses
(\cref{fig:overview}).

In this paper we propose \textbf{GAHIB}, which integrates three
synergistic inductive biases into a single VAE:
\begin{enumerate}[leftmargin=*,itemsep=2pt]
\item \textbf{Graph Attention Network (GAT) encoder.}
  A two-layer GAT operates on a $k$-nearest-neighbour cell graph,
  learning attention-weighted neighbourhood aggregation that captures
  cell-type-aware interactions.
\item \textbf{Information Bottleneck (IB) layer.}
  A linear compression maps the $d$-dimensional latent code $\bz$ to a
  2D manifold coordinate~$\bz'$, creating a structured low-dimensional
  representation amenable to geometric losses.
\item \textbf{Lorentz hyperbolic geometry loss.}
  Geodesic distance on the Lorentz hyperboloid between $\bz$ and $\bz'$
  encourages the latent space to develop hierarchical organisation,
  placing stem cells near the origin and differentiated cells at higher
  norms.
\end{enumerate}

A key architectural feature is that the IB layer supplies the compressed
manifold coordinate~$\bz'$ required to define a stable target for the
hyperbolic loss, while the GAT encoder strengthens both geometric terms
through neighbourhood-aware representations.

We validate GAHIB through eleven complementary experimental studies
organised into three groups.
Seven \emph{comparative benchmarks} quantify where GAHIB stands against
prior art:
\begin{enumerate}[leftmargin=*,itemsep=0pt,label=(\roman*)]
\item Component ablation (5~variants, \cref{sec:ablation});
\item Deep learning benchmark (8~methods, \cref{sec:benchmark_dl});
\item Classical DR benchmark (6~methods, \cref{sec:benchmark_classical});
\item Geometric VAE benchmark (6~methods, \cref{sec:benchmark_gmvae});
\item Disentanglement comparison (6~methods, \cref{sec:benchmark_disent});
\item Encoder architecture comparison (3~types, \cref{sec:encoder});
\item Graph convolution sweep (6~operators, \cref{sec:gconv}).
\end{enumerate}
A \emph{biological interpretation} study
(\cref{sec:interpretation}) asks whether the learned latent space
recovers developmental hierarchy, gene programs, cell-type homophily,
and tissue-specific Gene Ontology Biological Process enrichment
(\cref{sec:go}).
Four \emph{robustness and efficiency} studies
(\cref{sec:robustness}) further quantify latent-dimension stability
(\cref{sec:latdim}), multi-seed reproducibility (\cref{sec:seeds}),
computational cost and scaling (\cref{sec:cost}), and hyperparameter
sensitivity (\cref{sec:hpsens}), directly addressing concerns raised
during peer review of deep-learning models for bioinformatics.
Each experiment is paired with mean$\pm$std performance estimates and
Wilcoxon signed-rank tests against GAHIB, enabling consistent statistical
comparison across study types.


% ============================================================================
\section{Related Work}\label{sec:related}
% ============================================================================

\subsection{Deep Generative Models for scRNA-seq}

scVI~\cite{lopez2018deep} introduced the VAE framework to single-cell
analysis with a negative-binomial decoder.
Subsequent methods extended this with contrastive learning
(CLEAR~\cite{han2022clear}), deep adaptive clustering
(scDAC~\cite{an2024scdac}), batch integration
(SCALEX~\cite{xiong2022scalex}), and graph-based encoders
(scGNN~\cite{wang2021scgnn}, scGCC~\cite{tian2023scgcc}).
scDHMap~\cite{tian2023scdhmap} uses a hyperbolic decoder to preserve
hierarchical distances, but employs an MLP encoder and does not incorporate
an information bottleneck.

\subsection{Hyperbolic Representations}

Hyperbolic spaces embed tree-like hierarchies with exponentially lower
distortion than Euclidean space~\cite{nickel2017poincare}.
Ganea \emph{et al.}~\cite{ganea2018hyperbolic} extended this to hyperbolic
neural networks operating on the Lorentz hyperboloid via the exponential
map---the geometric foundation we adopt in GAHIB.
GM-VAE~\cite{bayer2023gmvae} explored five hyperbolic prior distributions
(Poincar\'e, PGM, Learnable PGM, Hyperbolic-Wasserstein) for single-cell
data, but relied on MLP encoders and did not jointly optimise manifold
geometry with an information bottleneck.
GAHIB advances this by coupling a GAT encoder with an explicit
IB$\to$hyperbolic geometry pipeline, enabling each component to reinforce
the others.

\subsection{Graph Neural Networks for Single-Cell Data}

Graph attention networks~\cite{velickovic2018graph} learn
attention-weighted message passing.
CellVGAE~\cite{buterez2022cellvgae} applied GCN encoders for single-cell
VAEs.
GAHIB extends this by combining GAT encoding with an information bottleneck,
enabling explicit geometric losses on the latent manifold.

\subsection{Information Bottleneck and Disentanglement}

The information bottleneck principle~\cite{tishby2000information}
compresses representations to retain task-relevant information while
discarding noise.
Alemi \emph{et al.}~\cite{alemi2017deep} cast this as a variational
objective (VIB), enabling end-to-end learning with reparameterised
stochastic encoders.
Disentanglement methods such as $\beta$-VAE~\cite{higgins2017betavae},
DIP-VAE~\cite{kumar2018dipvae}, $\beta$-TC-VAE~\cite{chen2018isolating},
FactorVAE~\cite{kim2018factor}, and InfoVAE~\cite{zhao2019infovae}
apply posterior regularisation to encourage statistically independent
latent dimensions.
Our empirical results show that these regularisers \emph{decrease}
clustering quality on scRNA-seq data, likely because they sacrifice
reconstruction fidelity for disentanglement.
GAHIB instead uses the IB layer as a \emph{geometric} bottleneck---a
structured anchor for hyperbolic distance losses---rather than a
disentanglement objective.


% ============================================================================
\section{Method}\label{sec:method}
% ============================================================================

\subsection{Problem Setting}

Given a scRNA-seq count matrix $\mathbf{X} \in \rea^{N \times G}$ with
$N$ cells and $G$ genes, we aim to learn a latent representation
$\mathbf{Z} \in \rea^{N \times d}$ that preserves biological structure
(cell types, trajectories, hierarchies) while discarding technical noise.

\subsection{Architecture Overview}

GAHIB is a variational autoencoder comprising three integrated modules
(\cref{fig:architecture}).

\begin{figure}[!htbp]
\centering
\includegraphics[width=\textwidth]{architecture}
% Source: GAHIB_results/figures/fig_architecture.pdf
% Generator: gahib/viz/fig_architecture.py
\caption{GAHIB architecture block diagram.
Preprocessed counts are converted to a $k$-NN graph, encoded by a
two-layer GAT, and sampled into a 10-dimensional latent code.
The latent representation is then constrained by an information
bottleneck, regularised in Lorentz-hyperbolic space, and decoded with a
negative-binomial head. The objective combines reconstruction, KL,
bottleneck, and hyperbolic losses.}
\label{fig:architecture}
\end{figure}

\paragraph{GAT Encoder.}
We construct a $k$-nearest-neighbour graph ($k{=}15$) over cells in
gene-expression space.
A two-layer GAT encoder maps each cell $\bx_i$ to variational
parameters $(\bmu_i, \log\bsig_i^2)$ via attention-weighted
neighbourhood aggregation:
\begin{equation}
  \bmu_i, \log\bsig_i^2 = \text{GAT}\bigl(\bx_i,
    \{\bx_j\}_{j \in \mathcal{N}(i)}\bigr)
\end{equation}
The latent code is sampled via the reparameterisation trick:
$\bz_i = \bmu_i + \bsig_i \odot \boldsymbol{\epsilon}$,
$\boldsymbol{\epsilon} \sim \mathcal{N}(\mathbf{0}, \mathbf{I})$.

\paragraph{Information Bottleneck.}
A linear projection compresses $\bz \in \rea^d$ to a 2D manifold
coordinate $\bz' \in \rea^2$, which is then reconstructed back to
$\hat{\bz} \in \rea^d$.
The IB reconstruction loss
$\mathcal{L}_\text{ib} = \lVert \bz - \hat{\bz} \rVert^2$
ensures information preservation through the bottleneck.
Crucially, $\bz'$ provides the \emph{anchor} for the hyperbolic
geometry loss---without IB, the manifold coordinate is untrained and
the Lorentz distance becomes degenerate.

\paragraph{Lorentz Hyperbolic Geometry.}
Both $\bz$ and $\bz'$ are mapped to the Lorentz hyperboloid
$\lore^d = \{x \in \rea^{d+1} : \langle x, x \rangle_\lore = -1,
x_0 > 0\}$
via the exponential map at the origin.
The hyperbolic geometry loss:
\begin{equation}
  \mathcal{L}_\text{hyp} = d_\lore\!\bigl(
    \text{ExpMap}(\bz),\;\text{ExpMap}(\bz')\bigr)
  \label{eq:hyp_loss}
\end{equation}
encourages hierarchical organisation by penalising geodesic distance
on the hyperboloid between the full and compressed representations.
The Lorentz norm $\lVert \text{ExpMap}(\bz_i) \rVert_\lore$ provides
a natural \emph{hierarchy score}: cells with lower norms occupy
positions closer to the hyperboloid origin (stem/progenitor), while
those with higher norms are pushed towards the boundary
(differentiated).

\paragraph{NB Decoder.}
An MLP decoder parameterises a negative-binomial likelihood
$p(\bx_i \mid \bz_i)$, capturing the overdispersion and
library-size variation of count data.

\subsection{Training Objective}

The total loss combines four terms:
\begin{equation}
  \mathcal{L} = \underbrace{\mathcal{L}_\text{recon}}_\text{NB likelihood}
              + \beta \cdot \underbrace{\KL\bigl(q(\bz|\bx)
                \| p(\bz)\bigr)}_\text{regularisation}
              + \lambda_\text{ib} \cdot
                \underbrace{\mathcal{L}_\text{ib}}_\text{bottleneck}
              + \lambda_\text{hyp} \cdot
                \underbrace{\mathcal{L}_\text{hyp}}_\text{geometry}
  \label{eq:total_loss}
\end{equation}
with default weights $\beta{=}0.1$, $\lambda_\text{ib}{=}0.5$,
$\lambda_\text{hyp}{=}5.0$.
The low $\beta$ permits flexible posterior approximation, while the
strong $\lambda_\text{hyp}$ anchors hierarchical structure.

\subsection{Subgraph Sampling}

Training the GAT encoder on the full cell graph is memory-intensive.
We employ \emph{subgraph sampling}: each epoch draws 10 subgraphs of
512 nodes from the full graph, providing sufficient gradient signal
to match the number of mini-batch updates for MLP encoders.
This was a critical implementation fix---without it, the graph encoder
received only one gradient step per epoch versus ${\sim}16$ for MLP,
leading to severe under-training.


\subsection{Post-Hoc Interpretation Analyses}

We define four post-hoc analyses that extract biological insight from
a trained GAHIB model without retraining.

\paragraph{Poincar\'e Projection and Lorentz Norm.}
Each latent code $\bz_i$ is mapped to the Lorentz hyperboloid via the
exponential map, yielding $\mathbf{h}_i = \text{ExpMap}_0(\bz_i)
\in \lore^{d+1}$.
The \emph{Lorentz norm} $r_i = \cosh^{-1}(h_{i,0})$ measures geodesic
distance from the hyperboloid origin.
We project to the Poincar\'e disk via the diffeomorphism
$\mathbf{p}_i = \mathbf{h}_{i,1:d+1} / (1 + h_{i,0})$,
providing a 2D visualisation where the radial coordinate encodes
hierarchical depth.

\paragraph{Stemness Score.}
We define a per-cell stemness score as the normalised inverse Lorentz
norm: $s_i = 1 - (r_i - r_\text{min}) / (r_\text{max} - r_\text{min})$.
The Spearman correlation $\rho_s$ between $s_i$ and external
developmental annotations (when available) or Leiden cluster ordering
quantifies the agreement between hyperbolic geometry and biological
hierarchy.

\paragraph{Decoder Jacobian Gene Attribution.}
For each cell $i$ we compute the decoder Jacobian
$\mathbf{J}_i = \partial g(\bz_i) / \partial \bz_i \in \rea^{G \times d}$,
where $g$ is the NB decoder mean function.
Aggregating $|\mathbf{J}_i|$ across a random subsample of 200 cells
yields a gene$\times$dimension importance matrix.
Genes with high column norms in this matrix are the top attributed
genes per latent dimension, representing learned gene programs.
Marker overlap (MO) is computed as the fraction of differentially
expressed marker genes (Wilcoxon rank-sum, Bonferroni-corrected
$p < 0.05$) that appear among the top-100 Jacobian-attributed genes.

\paragraph{Lorentz-Norm Pseudotime.}
We use the Lorentz norm itself as an unsupervised pseudotime:
$r_i = \cosh^{-1}(h_{i,0})$ where $\mathbf{h}_i = \text{ExpMap}_0(\bz_i)$.
The cell with the smallest $r_i$ is taken as the developmental root.
We validate $r_i$ against scanpy diffusion pseudotime (DPT,
computed from PCA$\to$$k$NN$\to$diffusion map with $k{=}15$) via
Spearman and Kendall rank correlations.
Two Euclidean baselines are included for comparison:
Euclidean distance in 30-D PCA space and in 2-D UMAP space, both
from the same root cell.

\paragraph{GO Biological Process Enrichment.}
For each latent dimension $L_k$ we compute the Pearson correlation
between $L_k$ and every gene across all cells, rank genes by
$|\rho|$, and take the top 100.
We query these gene lists against the MSigDB GO~BP~v2024.1 gene sets
(organism-matched: Mouse Mm for mouse datasets, Human Hs for human)
using \texttt{gseapy.enrich} with local GMT files.
Terms are filtered at Benjamini--Hochberg adjusted $p<0.05$.
For the cross-dataset summary we take the minimum adjusted $p$-value
per term across datasets and report the top three most significant
terms per latent dimension.

\paragraph{GAT Attention Homophily.}
We extract the learned attention coefficients $\alpha_{ij}$ from the
final GAT layer for all edges $(i,j)$ in the cell graph.
Attention homophily is defined as the attention-weighted fraction of
same-type neighbours:
$H = \sum_{(i,j)} \alpha_{ij} \cdot \mathbb{1}[y_i = y_j]
     \,/\, \sum_{(i,j)} \alpha_{ij}$,
where $y_i$ is the Leiden cluster label of cell~$i$.
Values near 1 indicate that the GAT preferentially attends to
biologically similar cells.


% ============================================================================
\section{Experimental Setup}\label{sec:experiments}
% ============================================================================

\subsection{Datasets}

\begin{sloppypar}
We evaluate on 53 scRNA-seq datasets spanning two biological domains
(\cref{fig:taxonomy}):
\begin{itemize}[leftmargin=*,itemsep=0pt]
\item \textbf{Cancer} (27): skin (basal and squamous cell carcinoma),
  breast (6~datasets including epithelial, stromal, and metastatic),
  liver (hepatocellular, hepatoblastoma, metastasis),
  blood/leukaemia (AML, ALL, multiple myeloma, lymphoma),
  GI tract (gastric, colorectal), lung adenocarcinoma (2),
  brain metastasis (2), Merkel cell carcinoma (2), and T-cell cancers.
\item \textbf{Development} (26): hematopoiesis (8~datasets from CD34+
  progenitors to aged HSCs), neural development (dentate gyrus, spinal
  cord, retina, astrocytes), embryonic/stem cell systems (hESC time
  series, endoderm), organ development (lung, pituitary, teeth),
  disease models (Alzheimer's, LPS response), and atlas-scale references
  (PanSci muscle, T-cell).
\end{itemize}
Datasets span both human and mouse, range from 2{,}531 to 957{,}975
cells pre-subsampling, and include both primary tumours and
developmental trajectories---providing a comprehensive test of GAHIB's
generalisability across biological contexts.

\begin{figure}[!htbp]
\centering
\includegraphics[width=\textwidth]{dataset_taxonomy}
% Source: GAHIB_results/figures/fig_dataset_taxonomy.pdf
% Generator: gahib/viz/fig_dataset_taxonomy.py
\caption{Sunburst taxonomy of the 53 scRNA-seq evaluation datasets.
The concentric rings summarise biological domain, subcategory, and
dataset identity, and the legend reports the corresponding counts.}
\label{fig:taxonomy}
\end{figure}

All datasets are preprocessed identically: library-size normalisation,
$\log(1{+}x)$ transform, selection of 2{,}000 highly variable genes,
and subsampling to at most 3{,}000 cells.
Cell-type labels are obtained via Leiden
clustering~\cite{traag2019leiden} (resolution\,$=$\,1.0) on
PCA$\to$UMAP neighbours computed with
Scanpy~\cite{wolf2018scanpy}, providing a fully unsupervised reference
throughout.
The random seed is fixed at~42 for reproducibility.
\end{sloppypar}

\subsection{Evaluation Metrics}

We assess latent representations along three axes comprising 20 metrics
total.

\textbf{Clustering quality} (6 metrics).
$K$-means clustering ($k$ matching the number of Leiden clusters) is
applied to the latent space, and the predicted assignments are compared
against the Leiden reference partition.
Normalised mutual information (NMI) and adjusted Rand index (ARI)
measure partition agreement; average silhouette width
(ASW$\uparrow$) quantifies intra- vs.\ inter-cluster distance;
Davies--Bouldin index (DAV$\downarrow$) penalises cluster overlap;
Calinski--Harabasz score (CAL$\uparrow$) rewards compact, well-separated
clusters; and mean absolute inter-dimensional Pearson correlation (COR)
captures redundancy among latent dimensions.

\textbf{Dimensionality reduction evaluation} (DRE, 8 metrics).
We evaluate neighbourhood preservation between the latent space and its
UMAP and t-SNE projections using the co-ranking framework with
$k{=}15$ neighbours.
For each projection we report distance correlation (rank-distance
agreement), $Q_\text{local}$ (local neighbourhood preservation),
$Q_\text{global}$ (global structure preservation), and an overall quality
score combining local and global components.

\textbf{Latent space evaluation} (LSE, 6 metrics).
We assess the intrinsic geometry of the latent space via
manifold dimensionality (intrinsic dimension estimate via PCA eigenvalue
spectrum), spectral decay rate (slope of the sorted eigenvalue curve),
participation ratio (effective number of active dimensions),
anisotropy score (directional spread uniformity), noise resilience
(stability under Gaussian perturbation), and an overall quality score
averaging the normalised component scores.

\subsection{Statistical Testing}

All pairwise comparisons use the \textbf{Wilcoxon signed-rank test}
(two-sided) across the 53 datasets, with significance levels:
\sigA~$p < 0.001$, \sigB~$p < 0.01$, \sigC~$p < 0.05$, and
\ns~not significant.
Tables report mean $\pm$ std with significance markers for each
method versus GAHIB (or GAT, in architecture-specific comparisons).
This non-parametric test is chosen because metric distributions across
heterogeneous datasets are not guaranteed to be normal, and the
paired design (same dataset, different methods) maximises statistical
power.

\subsection{Training Configuration}

All models are trained for 200 epochs with patience$=$30 (early stopping),
hidden dimension$=$128, latent dimension$=$10, IB dimension$=$2,
learning rate$=10^{-4}$, batch size$=$128, and an NB likelihood.
GAHIB-specific weights: $\beta{=}0.1$, $\lambda_\text{ib}{=}0.5$,
$\lambda_\text{hyp}{=}5.0$, with a GAT encoder on a 15-NN graph and
10$\times$512-node subgraph sampling.


% ============================================================================
\section{Results}\label{sec:results}
% ============================================================================

% ── 5.1 Ablation ─────────────────────────────────────────────────────────
\subsection{Component Ablation}\label{sec:ablation}

We ablate GAHIB by progressively removing components, yielding five
variants: Base VAE (MLP, no IB, no geometry), VAE+IB, VAE+Hyp,
VAE+IB+Hyp, and the full GAHIB (GAT+IB+Hyp).

\Cref{tab:ablation} reports mean$\pm$std across 53 datasets with Wilcoxon
significance against GAHIB.
\Cref{fig:ablation} shows boxplots across all 20~metrics.

\textbf{Information Bottleneck (+IB).}
Adding IB to the Base VAE improves ARI by $+$0.027, ASW by $+$0.044, and
reduces DAV by $-$0.252.
Mechanistically, the bottleneck forces the encoder to discard technical
noise and retain only the most discriminative gene-expression patterns,
which directly improves cluster separation (higher ASW, lower DAV).
The IB also provides the structured 2D manifold coordinate~$\bz'$
required for the hyperbolic geometry loss in subsequent variants.

\textbf{Hyperbolic Geometry (+Hyp).}
Adding Lorentz geometry alone (VAE+Hyp) yields only marginal clustering
improvement (NMI $+$0.002) but substantially boosts DRE
($0.553 \to 0.615$) and LSE ($0.194 \to 0.297$), indicating that the
hyperbolic loss regularises the \emph{global structure} of the latent
manifold even when cluster separation is unchanged.
Crucially, \emph{Hyp without IB is ineffective for clustering}: without a
trained manifold coordinate~$\bz'$, the Lorentz distance is computed
against an untrained target and degenerates.

\textbf{Combined IB+Hyp.}
VAE+IB+Hyp achieves the best DRE (0.735) and LSE (0.534) among all MLP
variants, confirming the synergy: IB provides the geometric anchor
enabling the hyperbolic loss to structure the full latent space, not
merely the 2D projection.
The $+$0.052 DRE gain over VAE+IB alone reflects improved global
neighbourhood preservation---a direct benefit of the Lorentz geodesic
penalty.

\textbf{Full GAHIB (GAT+IB+Hyp).}
Replacing the MLP encoder with GAT yields the largest single improvement:
NMI $+$0.072, ARI $+$0.078, ASW $+$0.149, DAV $-$0.425, all significant
at $p < 0.001$.
This gap arises because the GAT encoder aggregates information from
cell-type-homophilic neighbours (mean attention homophily 0.812),
producing pre-clustered representations that make the IB and hyperbolic
losses substantially more effective.
The GAT encoder is the \emph{critical enabler} for GAHIB's clustering
performance; the geometric losses then refine the latent geometry.

% --- Ablation Table with Significance ---
\begin{table}[!htbp]
\centering
\caption{Component ablation: mean $\pm$ std across 53 datasets.
Entries are compared with GAHIB using Wilcoxon signed-rank tests;
DAV$\downarrow$ indicates that lower values are better.}
\label{tab:ablation}
\resizebox{\textwidth}{!}{%
\begin{tabular}{l c c c c c c}
\toprule
Method & NMI$\uparrow$ & ARI$\uparrow$ & ASW$\uparrow$ & DAV$\downarrow$ & DRE$\uparrow$ & LSE$\uparrow$ \\
\midrule
Base VAE\sigA
  & $0.651{\scriptstyle\pm.141}$
  & $0.485{\scriptstyle\pm.152}$
  & $0.165{\scriptstyle\pm.038}$
  & $1.741{\scriptstyle\pm.160}$
  & $0.553{\scriptstyle\pm.039}$
  & $0.194{\scriptstyle\pm.025}$ \\
VAE+IB\sigA
  & $0.665{\scriptstyle\pm.137}$
  & $0.512{\scriptstyle\pm.154}$
  & $0.209{\scriptstyle\pm.051}$
  & $1.489{\scriptstyle\pm.185}$
  & $0.683{\scriptstyle\pm.042}$
  & $0.378{\scriptstyle\pm.041}$ \\
VAE+Hyp\sigA
  & $0.653{\scriptstyle\pm.140}$
  & $0.493{\scriptstyle\pm.153}$
  & $0.168{\scriptstyle\pm.041}$
  & $1.679{\scriptstyle\pm.168}$
  & $0.615{\scriptstyle\pm.048}$
  & $0.297{\scriptstyle\pm.048}$ \\
VAE+IB+Hyp\sigA
  & $0.650{\scriptstyle\pm.133}$
  & $0.486{\scriptstyle\pm.150}$
  & $0.215{\scriptstyle\pm.048}$
  & $1.402{\scriptstyle\pm.166}$
  & $0.735{\scriptstyle\pm.051}$
  & $0.534{\scriptstyle\pm.081}$ \\
\textbf{GAHIB}
  & $\mathbf{0.722}{\scriptstyle\pm.098}$
  & $\mathbf{0.564}{\scriptstyle\pm.142}$
  & $\mathbf{0.364}{\scriptstyle\pm.081}$
  & $\mathbf{0.977}{\scriptstyle\pm.139}$
  & $\mathbf{0.699}{\scriptstyle\pm.059}$
  & $\mathbf{0.576}{\scriptstyle\pm.086}$ \\
\bottomrule
\end{tabular}%
}
\end{table}

\begin{figure}[!htbp]
\centering
\includegraphics[width=\textwidth]{ablation_all_metrics}
% Source: GAHIB_results/ablation/figures/fig_all_metrics.pdf
% Generator: gahib/viz/controller.py via run_all_visualizations.py
\caption{Component ablation across all 20 evaluation metrics
across clustering, dimensionality-reduction, and latent-structure
metrics. Boxes show the interquartile range over 53 datasets, and
brackets denote Wilcoxon tests versus GAHIB.}
\label{fig:ablation}
\end{figure}


% ── 5.2 Deep Learning Benchmark ──────────────────────────────────────────
\subsection{Deep Learning Benchmark}\label{sec:benchmark_dl}

We compare GAHIB against seven published deep learning methods:
scVI~\cite{lopez2018deep}, CellBLAST~\cite{cao2020cellblast},
CLEAR~\cite{han2022clear}, SCALEX~\cite{xiong2022scalex},
scDeepCluster~\cite{tian2019scdeepcluster},
scDHMap~\cite{tian2023scdhmap}, and scGNN~\cite{wang2021scgnn}.
All are trained with 200~epochs and batch size 128 on identically
preprocessed data.
The paper-ready comparison retains the seven baselines that produced
stable, interpretable embeddings under the shared 53-dataset evaluation
pipeline.

\Cref{tab:scdeep} reports performance with significance against GAHIB.
GAHIB ranks first on all six aggregate metrics, significantly
outperforming all methods except scVI on NMI ($p = 0.059$, ns) and
ARI ($p = 0.989$, ns).
The runner-up, scDHMap, achieves moderate ASW (0.335) and DRE (0.711)
but lags substantially on clustering (NMI 0.612 vs.~0.715).
GAHIB provides the most \emph{balanced} profile across all metric
families (\cref{fig:scdeep}).

% --- SC Deep Learning Table with Significance ---
\begin{table}[!htbp]
\centering
\caption{Deep learning benchmark: mean $\pm$ std across 53 datasets.
Entries are compared with GAHIB using Wilcoxon signed-rank tests;
DAV$\downarrow$ indicates that lower values are better.}
\label{tab:scdeep}
\resizebox{\textwidth}{!}{%
\begin{tabular}{l c c c c c c}
\toprule
Method & NMI$\uparrow$ & ARI$\uparrow$ & ASW$\uparrow$ & DAV$\downarrow$ & DRE$\uparrow$ & LSE$\uparrow$ \\
\midrule
scVI\ns
  & $0.706{\scriptstyle\pm.116}$
  & $0.555{\scriptstyle\pm.141}$
  & $0.174{\scriptstyle\pm.039}$
  & $1.699{\scriptstyle\pm.162}$
  & $0.542{\scriptstyle\pm.037}$
  & $0.180{\scriptstyle\pm.015}$ \\
CellBLAST\sigA
  & $0.049{\scriptstyle\pm.023}$
  & $0.017{\scriptstyle\pm.014}$
  & $0.026{\scriptstyle\pm.014}$
  & $1.993{\scriptstyle\pm.174}$
  & $0.424{\scriptstyle\pm.032}$
  & $0.157{\scriptstyle\pm.014}$ \\
CLEAR\sigA
  & $0.506{\scriptstyle\pm.126}$
  & $0.338{\scriptstyle\pm.119}$
  & $0.118{\scriptstyle\pm.021}$
  & $1.917{\scriptstyle\pm.108}$
  & $0.421{\scriptstyle\pm.024}$
  & $0.132{\scriptstyle\pm.009}$ \\
SCALEX\sigA
  & $0.045{\scriptstyle\pm.027}$
  & $0.016{\scriptstyle\pm.022}$
  & $0.045{\scriptstyle\pm.086}$
  & $1.998{\scriptstyle\pm.171}$
  & $0.418{\scriptstyle\pm.028}$
  & $0.158{\scriptstyle\pm.014}$ \\
scDeepCluster\sigA
  & $0.574{\scriptstyle\pm.096}$
  & $0.380{\scriptstyle\pm.104}$
  & $0.129{\scriptstyle\pm.028}$
  & $1.813{\scriptstyle\pm.140}$
  & $0.497{\scriptstyle\pm.051}$
  & $0.279{\scriptstyle\pm.051}$ \\
scDHMap\sigA
  & $0.612{\scriptstyle\pm.111}$
  & $0.433{\scriptstyle\pm.122}$
  & $0.335{\scriptstyle\pm.061}$
  & $1.032{\scriptstyle\pm.153}$
  & $0.711{\scriptstyle\pm.057}$
  & $0.518{\scriptstyle\pm.114}$ \\
scGNN\sigA
  & $0.173{\scriptstyle\pm.049}$
  & $0.068{\scriptstyle\pm.026}$
  & $0.122{\scriptstyle\pm.018}$
  & $1.725{\scriptstyle\pm.110}$
  & $0.525{\scriptstyle\pm.019}$
  & $0.229{\scriptstyle\pm.027}$ \\
\textbf{GAHIB}
  & $\mathbf{0.715}{\scriptstyle\pm.117}$
  & $\mathbf{0.560}{\scriptstyle\pm.167}$
  & $\mathbf{0.359}{\scriptstyle\pm.087}$
  & $\mathbf{1.002}{\scriptstyle\pm.186}$
  & $\mathbf{0.704}{\scriptstyle\pm.057}$
  & $\mathbf{0.572}{\scriptstyle\pm.100}$ \\
\bottomrule
\end{tabular}%
}
\end{table}

\begin{figure}[!htbp]
\centering
\includegraphics[width=\textwidth]{scdeep_all_metrics}
% Source: GAHIB_results/sc_deeplearning_benchmark/figures/fig_all_metrics.pdf
% Generator: gahib/viz/controller.py via run_all_visualizations.py
\caption{Deep learning benchmark across all 20 metrics for 8 methods.
GAHIB leads the aggregate metrics and remains
statistically superior to most baselines; detailed tests are summarised
in \cref{tab:scdeep}.}
\label{fig:scdeep}
\end{figure}


% ── 5.3 Classical DR Benchmark ───────────────────────────────────────────
\subsection{Classical Dimensionality Reduction Benchmark}
\label{sec:benchmark_classical}

We compare GAHIB with PCA, ICA, NMF, Truncated SVD, and Diffusion Maps,
each extracting 10 components from the same preprocessed data.

\Cref{tab:classical} shows that GAHIB achieves the best DRE (0.707) and
LSE (0.576), with significance at $p < 0.001$ versus all classical
methods on these families.
PCA and Truncated SVD achieve the highest NMI (0.735, 0.732) and ARI (0.586, 0.584),
reflecting strong performance on linearly separable data.
DiffusionMap leads on ASW (0.563) and DAV (0.645), benefiting from its
nonlinear manifold structure; however, its clustering quality is
substantially lower (NMI 0.693).
GAHIB provides the best overall balance, dominating on complex
datasets where nonlinear structure matters (\cref{fig:classical}).

% --- Classical Benchmark Table with Significance ---
\begin{table}[!htbp]
\centering
\caption{Classical dimensionality-reduction benchmark: mean $\pm$ std
across 53 datasets. Entries are compared with GAHIB using Wilcoxon
signed-rank tests; DAV$\downarrow$ indicates that lower values are better.}
\label{tab:classical}
\resizebox{\textwidth}{!}{%
\begin{tabular}{l c c c c c c}
\toprule
Method & NMI$\uparrow$ & ARI$\uparrow$ & ASW$\uparrow$ & DAV$\downarrow$ & DRE$\uparrow$ & LSE$\uparrow$ \\
\midrule
PCA\sigA
  & $\mathbf{0.735}{\scriptstyle\pm.075}$
  & $\mathbf{0.586}{\scriptstyle\pm.122}$
  & $0.275{\scriptstyle\pm.067}$
  & $1.285{\scriptstyle\pm.168}$
  & $0.676{\scriptstyle\pm.068}$
  & $0.336{\scriptstyle\pm.069}$ \\
ICA\sigA
  & $0.697{\scriptstyle\pm.097}$
  & $0.547{\scriptstyle\pm.131}$
  & $0.261{\scriptstyle\pm.067}$
  & $1.282{\scriptstyle\pm.184}$
  & $0.571{\scriptstyle\pm.070}$
  & $0.139{\scriptstyle\pm.060}$ \\
NMF\sigA
  & $0.691{\scriptstyle\pm.094}$
  & $0.524{\scriptstyle\pm.134}$
  & $0.310{\scriptstyle\pm.068}$
  & $1.148{\scriptstyle\pm.140}$
  & $0.659{\scriptstyle\pm.072}$
  & $0.294{\scriptstyle\pm.061}$ \\
Trunc.\ SVD\sigA
  & $0.732{\scriptstyle\pm.072}$
  & $0.584{\scriptstyle\pm.118}$
  & $0.271{\scriptstyle\pm.065}$
  & $1.299{\scriptstyle\pm.164}$
  & $0.678{\scriptstyle\pm.069}$
  & $0.339{\scriptstyle\pm.071}$ \\
DiffusionMap\sigA
  & $0.693{\scriptstyle\pm.097}$
  & $0.508{\scriptstyle\pm.151}$
  & $\mathbf{0.563}{\scriptstyle\pm.136}$
  & $\mathbf{0.645}{\scriptstyle\pm.165}$
  & $0.654{\scriptstyle\pm.043}$
  & $0.150{\scriptstyle\pm.018}$ \\
\textbf{GAHIB}
  & $0.721{\scriptstyle\pm.106}$
  & $0.570{\scriptstyle\pm.150}$
  & $0.362{\scriptstyle\pm.089}$
  & $0.991{\scriptstyle\pm.183}$
  & $\mathbf{0.707}{\scriptstyle\pm.058}$
  & $\mathbf{0.576}{\scriptstyle\pm.095}$ \\
\bottomrule
\end{tabular}%
}
\end{table}

\begin{figure}[!htbp]
\centering
\includegraphics[width=\textwidth]{classical_all_metrics}
% Source: GAHIB_results/classical_benchmark/figures/fig_all_metrics.pdf
% Generator: gahib/viz/controller.py via run_all_visualizations.py
\caption{Classical dimensionality-reduction benchmark across 20 metrics
for 6 methods. GAHIB leads the DRE and LSE families, while PCA and
truncated SVD remain competitive on linearly separable datasets.}
\label{fig:classical}
\end{figure}


% ── 5.4 Geometric VAE Benchmark ──────────────────────────────────────────
\subsection{Geometric VAE Benchmark}\label{sec:benchmark_gmvae}

We compare GAHIB against five GM-VAE~\cite{bayer2023gmvae} geometric
prior configurations: Euclidean, Poincar\'e, PGM, Learnable PGM (L-PGM),
and Hyperbolic-Wasserstein (HW).
All GM-VAE variants use MLP encoders.

\Cref{tab:gmvae} shows GAHIB outperforms all variants across every metric.
Only the Learnable PGM variant achieves competitive clustering
(NMI\,$=$\,0.600), while the remaining variants fail catastrophically
(NMI\,$\leq$\,0.027).
This highlights that hyperbolic priors alone are insufficient---the
combination of GAT encoding with IB-mediated geometry is essential
(\cref{fig:gmvae}).

% --- GM-VAE Table with Significance ---
\begin{table}[!htbp]
\centering
\caption{Geometric VAE benchmark: mean $\pm$ std across 53 datasets.
Entries are compared with GAHIB using Wilcoxon signed-rank tests.}
\label{tab:gmvae}
\resizebox{\textwidth}{!}{%
\begin{tabular}{l c c c c c c}
\toprule
Method & NMI$\uparrow$ & ARI$\uparrow$ & ASW$\uparrow$ & DAV$\downarrow$ & DRE$\uparrow$ & LSE$\uparrow$ \\
\midrule
GM-VAE (Eucl.)\sigA
  & $0.019{\scriptstyle\pm.013}$
  & $0.001{\scriptstyle\pm.004}$
  & $0.035{\scriptstyle\pm.001}$
  & $2.984{\scriptstyle\pm.127}$
  & $0.329{\scriptstyle\pm.006}$
  & $0.105{\scriptstyle\pm.009}$ \\
GM-VAE (Poinc.)\sigA
  & $0.027{\scriptstyle\pm.013}$
  & $0.003{\scriptstyle\pm.002}$
  & $0.077{\scriptstyle\pm.003}$
  & $2.036{\scriptstyle\pm.076}$
  & $0.402{\scriptstyle\pm.005}$
  & $0.127{\scriptstyle\pm.009}$ \\
GM-VAE (PGM)\sigA
  & $0.018{\scriptstyle\pm.011}$
  & $0.000{\scriptstyle\pm.001}$
  & $0.069{\scriptstyle\pm.003}$
  & $2.126{\scriptstyle\pm.072}$
  & $0.452{\scriptstyle\pm.009}$
  & $0.131{\scriptstyle\pm.004}$ \\
GM-VAE (L-PGM)\sigA
  & $0.600{\scriptstyle\pm.147}$
  & $0.404{\scriptstyle\pm.149}$
  & $0.182{\scriptstyle\pm.054}$
  & $1.608{\scriptstyle\pm.248}$
  & $0.715{\scriptstyle\pm.066}$
  & $0.415{\scriptstyle\pm.084}$ \\
GM-VAE (HW)\sigA
  & $0.025{\scriptstyle\pm.010}$
  & $0.001{\scriptstyle\pm.004}$
  & $0.181{\scriptstyle\pm.042}$
  & $1.185{\scriptstyle\pm.083}$
  & $0.555{\scriptstyle\pm.005}$
  & $0.156{\scriptstyle\pm.041}$ \\
\textbf{GAHIB}
  & $\mathbf{0.726}{\scriptstyle\pm.104}$
  & $\mathbf{0.577}{\scriptstyle\pm.154}$
  & $\mathbf{0.365}{\scriptstyle\pm.079}$
  & $\mathbf{0.973}{\scriptstyle\pm.143}$
  & $\mathbf{0.707}{\scriptstyle\pm.050}$
  & $\mathbf{0.562}{\scriptstyle\pm.078}$ \\
\bottomrule
\end{tabular}%
}
\end{table}

\begin{figure}[!htbp]
\centering
\includegraphics[width=\textwidth]{gmvae_all_metrics}
% Source: GAHIB_results/gmvae_benchmark/figures/fig_all_metrics.pdf
% Generator: gahib/viz/controller.py via run_all_visualizations.py
\caption{Geometric VAE benchmark across 20 metrics for 6 methods.
Only the learnable PGM variant remains competitive; the other geometric
priors collapse to weak embeddings. GAHIB outperforms all variants.}
\label{fig:gmvae}
\end{figure}


% ── 5.5 Disentanglement Comparison ───────────────────────────────────────
\subsection{Disentanglement Regularisation Comparison}
\label{sec:benchmark_disent}

We compare GAHIB against four disentanglement regularisation strategies
applied to a base VAE with MLP encoder: $\beta$-VAE, DIP-VAE,
$\beta$-TC-VAE, and InfoVAE.

\Cref{tab:disent} shows that GAHIB outperforms all methods with
$p < 0.001$ significance across every metric.
Strikingly, $\beta$-VAE, DIP-VAE, and TC-VAE \emph{decrease} clustering
quality below the unregularised Base VAE (TC-VAE NMI 0.404 vs.\ Base
0.647).
This result has a clear mechanistic explanation: these regularisers impose
strong constraints on the posterior distribution---encouraging independent
or factorised latent dimensions---which disrupts the encoder's ability to
form coherent cluster representations.
In scRNA-seq data, meaningful biological structure \emph{requires}
correlated latent dimensions (e.g., co-expressed gene programs), so
disentanglement regularisation actively harms the objective.
InfoVAE, which relaxes the factorisation constraint, is nearly neutral
(NMI 0.649 vs.\ Base 0.647), consistent with this interpretation.
GAHIB achieves NMI 0.722 ($+$0.073 over InfoVAE, $+$0.318 over TC-VAE),
demonstrating that a \emph{geometric} structural prior (IB+Hyp) is
categorically more effective than posterior regularisation for
single-cell clustering (\cref{fig:disent}).

% --- Disentanglement Table with Significance ---
\begin{table}[!htbp]
\centering
\caption{Disentanglement comparison: mean $\pm$ std across 53 datasets.
Entries are compared with GAHIB using Wilcoxon signed-rank tests.}
\label{tab:disent}
\resizebox{\textwidth}{!}{%
\begin{tabular}{l c c c c c c}
\toprule
Method & NMI$\uparrow$ & ARI$\uparrow$ & ASW$\uparrow$ & DAV$\downarrow$ & DRE$\uparrow$ & LSE$\uparrow$ \\
\midrule
Base VAE\sigA
  & $0.647{\scriptstyle\pm.138}$
  & $0.480{\scriptstyle\pm.149}$
  & $0.164{\scriptstyle\pm.039}$
  & $1.740{\scriptstyle\pm.155}$
  & $0.548{\scriptstyle\pm.041}$
  & $0.195{\scriptstyle\pm.026}$ \\
$\beta$-VAE\sigA
  & $0.513{\scriptstyle\pm.156}$
  & $0.337{\scriptstyle\pm.151}$
  & $0.119{\scriptstyle\pm.030}$
  & $1.980{\scriptstyle\pm.155}$
  & $0.495{\scriptstyle\pm.036}$
  & $0.161{\scriptstyle\pm.020}$ \\
DIP-VAE\sigA
  & $0.491{\scriptstyle\pm.122}$
  & $0.306{\scriptstyle\pm.116}$
  & $0.105{\scriptstyle\pm.021}$
  & $2.142{\scriptstyle\pm.160}$
  & $0.467{\scriptstyle\pm.047}$
  & $0.159{\scriptstyle\pm.010}$ \\
TC-VAE\sigA
  & $0.404{\scriptstyle\pm.193}$
  & $0.234{\scriptstyle\pm.160}$
  & $0.109{\scriptstyle\pm.030}$
  & $1.936{\scriptstyle\pm.131}$
  & $0.551{\scriptstyle\pm.084}$
  & $0.263{\scriptstyle\pm.102}$ \\
InfoVAE\sigA
  & $0.649{\scriptstyle\pm.140}$
  & $0.484{\scriptstyle\pm.151}$
  & $0.164{\scriptstyle\pm.038}$
  & $1.754{\scriptstyle\pm.166}$
  & $0.551{\scriptstyle\pm.040}$
  & $0.190{\scriptstyle\pm.027}$ \\
\textbf{GAHIB}
  & $\mathbf{0.722}{\scriptstyle\pm.113}$
  & $\mathbf{0.576}{\scriptstyle\pm.161}$
  & $\mathbf{0.356}{\scriptstyle\pm.086}$
  & $\mathbf{0.994}{\scriptstyle\pm.192}$
  & $\mathbf{0.700}{\scriptstyle\pm.058}$
  & $\mathbf{0.559}{\scriptstyle\pm.093}$ \\
\bottomrule
\end{tabular}%
}
\end{table}

\begin{figure}[!htbp]
\centering
\includegraphics[width=\textwidth]{disent_all_metrics}
% Source: GAHIB_results/disentanglement/figures/fig_all_metrics.pdf
% Generator: gahib/viz/controller.py via run_all_visualizations.py
\caption{Disentanglement regularisation comparison across 20 metrics.
Posterior regularisers reduce clustering quality relative to the base VAE,
whereas GAHIB's structural prior remains consistently stronger.}
\label{fig:disent}
\end{figure}


% ── 5.6 Encoder Architecture ─────────────────────────────────────────────
\subsection{Encoder Architecture Comparison}\label{sec:encoder}

To isolate the effect of the encoder, we fix the GAHIB loss
configuration (IB+Hyp) and vary only the encoder: MLP, Transformer
(multi-head self-attention), and GAT.

\Cref{tab:encoder} shows a clear hierarchy: GAT $>$ Transformer $>$
MLP across all clustering metrics, all significant at $p < 0.001$.
GAT achieves NMI 0.723 vs.\ Transformer 0.659 and MLP 0.640.
The GAT encoder benefits from explicit graph structure (15-NN graph),
while the Transformer's self-attention provides partial neighbourhood
awareness.
This confirms that graph-structured encoding is a stronger inductive
bias than self-attention alone (\cref{fig:encoder}).

% --- Encoder Table with Significance ---
\begin{table}[!htbp]
\centering
\caption{Encoder comparison: mean $\pm$ std across 53 datasets.
Entries are compared with the GAT encoder using Wilcoxon signed-rank tests.}
\label{tab:encoder}
\resizebox{\textwidth}{!}{%
\begin{tabular}{l c c c c c c}
\toprule
Encoder & NMI$\uparrow$ & ARI$\uparrow$ & ASW$\uparrow$ & DAV$\downarrow$ & DRE$\uparrow$ & LSE$\uparrow$ \\
\midrule
MLP\sigA
  & $0.640{\scriptstyle\pm.132}$
  & $0.479{\scriptstyle\pm.153}$
  & $0.210{\scriptstyle\pm.057}$
  & $1.370{\scriptstyle\pm.193}$
  & $0.677{\scriptstyle\pm.054}$
  & $0.490{\scriptstyle\pm.091}$ \\
Transformer\sigA
  & $0.659{\scriptstyle\pm.129}$
  & $0.504{\scriptstyle\pm.158}$
  & $0.246{\scriptstyle\pm.059}$
  & $1.238{\scriptstyle\pm.150}$
  & $0.695{\scriptstyle\pm.056}$
  & $0.583{\scriptstyle\pm.121}$ \\
\textbf{GAT}
  & $\mathbf{0.723}{\scriptstyle\pm.107}$
  & $\mathbf{0.571}{\scriptstyle\pm.152}$
  & $\mathbf{0.359}{\scriptstyle\pm.083}$
  & $\mathbf{0.993}{\scriptstyle\pm.170}$
  & $\mathbf{0.709}{\scriptstyle\pm.052}$
  & $\mathbf{0.570}{\scriptstyle\pm.097}$ \\
\bottomrule
\end{tabular}%
}
\end{table}

\begin{figure}[!htbp]
\centering
\includegraphics[width=\textwidth]{encoder_all_metrics}
% Source: GAHIB_results/encoder_comparison/figures/fig_all_metrics.pdf
% Generator: gahib/viz/controller.py via run_all_visualizations.py
\caption{Encoder architecture comparison across 20 metrics with the
IB+Hyp objective fixed. GAT consistently outperforms Transformer and MLP,
showing that explicit graph structure is more effective than
self-attention alone.}
\label{fig:encoder}
\end{figure}


% ── 5.7 Graph Convolution Sweep ──────────────────────────────────────────
\subsection{Graph Convolution Operator Sweep}\label{sec:gconv}

We compare six graph convolution operators within the GAHIB framework:
GCN, GraphSAGE, Chebyshev (Cheb), TAG, GraphTransformer, and GAT.

\Cref{tab:gconv} shows GAT as the strongest operator overall.
GAT achieves the best NMI (0.744), ARI (0.551), ASW (0.358),
DAV (0.990), and LSE (0.591), while GCN offers a competitive
alternative (NMI 0.727).
The attention mechanism in GAT provides learnable, cell-specific
aggregation weights, which improves cluster separation
beyond fixed-weight message passing (\cref{fig:gconv}).

% --- Graph Conv Table with Significance ---
\begin{table}[!htbp]
\centering
\caption{Graph convolution sweep: mean $\pm$ std across 53 datasets.
Entries are compared with GAT using Wilcoxon signed-rank tests.}
\label{tab:gconv}
\resizebox{\textwidth}{!}{%
\begin{tabular}{l c c c c c c}
\toprule
Conv.\ Type & NMI$\uparrow$ & ARI$\uparrow$ & ASW$\uparrow$ & DAV$\downarrow$ & DRE$\uparrow$ & LSE$\uparrow$ \\
\midrule
GCN\sigA
  & $0.727{\scriptstyle\pm.104}$
  & $0.522{\scriptstyle\pm.150}$
  & $0.314{\scriptstyle\pm.073}$
  & $1.102{\scriptstyle\pm.167}$
  & $0.693{\scriptstyle\pm.065}$
  & $0.558{\scriptstyle\pm.074}$ \\
GraphSAGE\sigA
  & $0.666{\scriptstyle\pm.113}$
  & $0.470{\scriptstyle\pm.133}$
  & $0.236{\scriptstyle\pm.058}$
  & $1.379{\scriptstyle\pm.147}$
  & $0.682{\scriptstyle\pm.049}$
  & $0.511{\scriptstyle\pm.088}$ \\
Cheb\sigA
  & $0.633{\scriptstyle\pm.120}$
  & $0.428{\scriptstyle\pm.139}$
  & $0.197{\scriptstyle\pm.054}$
  & $1.548{\scriptstyle\pm.143}$
  & $0.670{\scriptstyle\pm.039}$
  & $0.445{\scriptstyle\pm.040}$ \\
TAG\sigA
  & $0.671{\scriptstyle\pm.115}$
  & $0.474{\scriptstyle\pm.142}$
  & $0.252{\scriptstyle\pm.060}$
  & $1.353{\scriptstyle\pm.139}$
  & $0.684{\scriptstyle\pm.054}$
  & $0.540{\scriptstyle\pm.069}$ \\
GraphTransformer\sigA
  & $0.661{\scriptstyle\pm.120}$
  & $0.460{\scriptstyle\pm.138}$
  & $0.237{\scriptstyle\pm.063}$
  & $1.361{\scriptstyle\pm.160}$
  & $0.692{\scriptstyle\pm.043}$
  & $0.519{\scriptstyle\pm.076}$ \\
\textbf{GAT}
  & $\mathbf{0.744}{\scriptstyle\pm.100}$
  & $\mathbf{0.551}{\scriptstyle\pm.152}$
  & $\mathbf{0.358}{\scriptstyle\pm.079}$
  & $\mathbf{0.990}{\scriptstyle\pm.158}$
  & $\mathbf{0.705}{\scriptstyle\pm.059}$
  & $\mathbf{0.591}{\scriptstyle\pm.091}$ \\
\bottomrule
\end{tabular}%
}
\end{table}

\begin{figure}[!htbp]
\centering
\includegraphics[width=\textwidth]{gconv_all_metrics}
% Source: GAHIB_results/graph_conv_sweep/figures/fig_all_metrics.pdf
% Generator: gahib/viz/controller.py via run_all_visualizations.py
\caption{Graph convolution operator sweep across 20 metrics within the
GAHIB framework. GAT is the strongest overall operator, while
GraphTransformer attains the highest DRE.}
\label{fig:gconv}
\end{figure}


% ============================================================================
\section{Biological Interpretation}\label{sec:interpretation}
% ============================================================================

We perform post-hoc analysis of GAHIB's trained representations across
all 53 datasets.
All results use the same trained models as the benchmark experiments.

\subsection{Latent Embedding and Poincar\'e Projection}

\Cref{fig:interp_embedding} shows UMAP embeddings (panel~a) and
Poincar\'e disk projections (panel~b) for all 53 datasets.
UMAP embeddings demonstrate clean cluster separation consistent with
GAHIB's high clustering scores.
The Poincar\'e projections reveal that cell types occupy distinct
regions of the hyperbolic disk, with differentiated types typically
located closer to the boundary (higher Lorentz norm), confirming
hierarchical organisation.

\begin{figure}[!htbp]
\centering
\includegraphics[width=\textwidth]{interp_embedding_poincare}
% Source: GAHIB_results/interpretation/figures/fig_themed_embedding_poincare.pdf
% Generator: gahib/viz/interpretation.py::fig_themed_embedding_poincare()
\caption{GAHIB latent representations across all 53 datasets.
\textbf{(a)}~UMAP embeddings coloured by Leiden cluster show clear
cluster separation.
\textbf{(b)}~Poincar\'e projections place differentiated cell states
closer to the boundary, consistent with hierarchical organisation.}
\label{fig:interp_embedding}
\end{figure}

\subsection{Information Bottleneck Structure}

\Cref{fig:interp_bottleneck} visualises the IB compression.
Panel~(a) shows 2D IB scatter coloured by cell type, demonstrating
that the bottleneck preserves cluster structure even at extreme
compression ($10\text{D} \to 2\text{D}$).
Panel~(b) overlays the most active latent dimension on UMAP,
showing that different dimensions capture distinct biological axes.
All 10 latent dimensions are active across all 53 datasets
(\cref{tab:interp}), with no dimension collapse.

\begin{figure}[!htbp]
\centering
\includegraphics[width=\textwidth]{interp_bottleneck}
% Source: GAHIB_results/interpretation/figures/fig_themed_bottleneck.pdf
% Generator: gahib/viz/interpretation.py::fig_themed_bottleneck()
\caption{Information bottleneck structure analysis.
\textbf{(a)}~2D IB scatter of the compressed manifold coordinate
$\bz' \in \rea^2$, coloured by Leiden cluster, showing that cluster
structure is retained after $10\text{D} \to 2\text{D}$ compression.
\textbf{(b)}~UMAP coloured by the most active latent dimension
(highest variance), indicating that individual dimensions capture
distinct biological axes.}
\label{fig:interp_bottleneck}
\end{figure}

\subsection{Hyperbolic Hierarchy}

\Cref{fig:interp_hyperbolic} colours cells by their Lorentz norm
on UMAP (panel~a) and Poincar\'e disk (panel~b).
The norm gradient is clearly visible: stem/progenitor cells have
lower norms (closer to origin) while differentiated cells have higher
norms (closer to boundary).

This is quantified by the stemness--norm Spearman correlation
$\rho_s$ in \cref{tab:interp}, which is significantly positive
in the majority of datasets (mean $\rho_s = 0.327$).
The strongest correlations appear in developmental datasets:
dentate ($\rho_s = 0.673$), breast cancer ($0.615$), MCCTumor
($0.501$), and colorectal adenocarcinoma ($0.505$).

\begin{figure}[!htbp]
\centering
\includegraphics[width=\textwidth]{interp_hyperbolic}
% Source: GAHIB_results/interpretation/figures/fig_themed_hyperbolic.pdf
% Generator: gahib/viz/interpretation.py::fig_themed_hyperbolic()
\caption{Hyperbolic hierarchy analysis.
\textbf{(a)}~UMAP embeddings coloured by Lorentz norm
$\lVert \text{ExpMap}(\bz_i) \rVert_\lore$; purple indicates low
norm (stem/progenitor), yellow indicates high norm (differentiated).
\textbf{(b)}~Poincar\'e disk projections coloured by the same norm,
showing the radial gradient from origin to boundary that encodes
developmental hierarchy.}
\label{fig:interp_hyperbolic}
\end{figure}

\subsection{Gene Ontology Biological Process Enrichment}\label{sec:go}

To further validate that individual GAHIB latent dimensions encode
coherent biological programmes, we performed GO Biological Process
enrichment analysis on four representative datasets (dentate gyrus,
hematopoiesis, Setty hematopoietic, hESC) covering neural,
haematopoietic, and embryonic lineages.
For each latent dimension $L_k$ we computed the Pearson correlation
between $L_k$ and every gene across cells, selected the top 100
correlated genes, and queried them against the MSigDB
GO~BP~v2024.1 gene sets (organism-matched) via \texttt{gseapy}.
Significance used Benjamini--Hochberg adjusted $p{<}0.05$.

Panels~(f)--(i) of \cref{fig:downstream} summarise the top-three
enriched terms per latent dimension for each tissue context.
The pattern reveals that GAHIB's latent space organises biological
programmes in a highly compartmentalised manner:
$L_0$ and $L_3$ capture immune-related pathways
(\textsc{Inflammatory Response}, \textsc{Defense Response};
$-\log_{10}$ adj.\,$p$ up to 22.7);
$L_1, L_4, L_5, L_9$ encode neural programmes
(\textsc{Neurogenesis}, \textsc{Neuron Development},
\textsc{CNS Development});
$L_2$ and $L_6$ recover cell-cycle machinery
(\textsc{Mitotic Cell Cycle}, \textsc{Chromosome Segregation});
and $L_7$ specifically captures glial biology
(\textsc{Glial Cell Differentiation}).
Haemostasis and platelet activation are localised to $L_8$ and $L_9$.
The absence of off-diagonal enrichment within each functional
category confirms that latent dimensions are functionally specialised
rather than redundant, and the high significance levels
(adj.\,$p\leq 10^{-9}$ throughout) confirm that the identified
pathways are not noise.

\paragraph{Tissue-context specificity.}
Because GAHIB is trained independently on each dataset, the latent
dimensions learn tissue-specific biological programmes rather than
generic pathways.
Panels~(f)--(i) of \cref{fig:downstream} display the top-three
enriched GO~BP terms per latent dimension for each of the four
datasets, trained as separate GAHIB models.
The enrichment is strikingly context-dependent:
\textbf{dentate gyrus} (mouse, neural) recovers
\textsc{Gliogenesis}, \textsc{Neurogenesis},
\textsc{Central Nervous System Development}, and
\textsc{Glial Cell Differentiation}
(adj.\,$p$ down to $10^{-21}$);
\textbf{hematopoiesis} (mouse, blood) surfaces
\textsc{Leukocyte Migration}, \textsc{Hemopoiesis},
\textsc{Granulocyte/Neutrophil Activation}, and
\textsc{Defense Response to Bacterium};
\textbf{Setty HSPC} (human, immune) yields
\textsc{B Cell Activation}, \textsc{Immune Effector Process},
and \textsc{Lymphocyte Activation};
\textbf{hESC} (human, stem/hemostatic) captures
\textsc{Hemostasis}, \textsc{Platelet Activation},
\textsc{Wound Healing}, and \textsc{Inflammatory Response}
(adj.\,$p$ down to $10^{-23}$).
Terms from one tissue context never appear in another---%
the model trained on brain tissue does not recover platelet biology,
and the blood-trained model does not recover neurogenesis.
This tissue specificity, discovered \emph{without any supervision},
indicates that GAHIB learns disentangled, biologically meaningful
programmes that reflect the underlying cellular biology rather than
generic transcriptional noise.

\subsection{Pseudotime Validation from Gene Expression}\label{sec:pseudotime}

Rather than evaluating the Lorentz-norm pseudotime against an
algorithmic reference (such as scanpy's diffusion pseudotime), we
validate it directly from the original gene expression: a faithful
developmental pseudotime should order cells so that known
progenitor-like programmes decrease along $r_i$ while
lineage-commitment programmes increase.
Concretely, for each dataset we (i)~train an independent GAHIB
model, (ii)~pick the cell with the smallest Lorentz norm as the
developmental root, (iii)~compute the Pearson correlation
$r_g = \mathrm{corr}(x_g, r_{\cdot})$ between every highly-variable
gene $x_g$ and GAHIB's pseudotime, and (iv)~report the top
negatively-correlated gene (early / progenitor-like) and the top
positively-correlated gene (lineage-commitment) per dataset.
The mean $|r|$ over the top-20 most-correlated HVGs summarises how
strongly GAHIB's ordering aligns with the dominant gene programme.

\Cref{fig:downstream} panels~(a)--(e) summarise the gene-expression
validation.  Across all four datasets the top negatively-correlated
and top positively-correlated HVGs show well-separated expression
trajectories along $r_i$ (panels~c,~d), and the mean $|r|$ of the
top-20 HVGs (panel~e) confirms that GAHIB's hyperbolic pseudotime
tracks a biologically coherent gene programme rather than a
statistical artefact of the embedding.

\subsection{Unified downstream analysis figure}\label{sec:downstream}

We consolidate the trajectory / pseudotime validation and the
tissue-specific GO enrichment into a single downstream-analysis
figure (\cref{fig:downstream}).
Tissue-context specificity is visible in panels~(f)--(i): four
GAHIB models, trained independently on the four datasets, recover
strikingly distinct GO~BP programmes with no cross-context leakage.
Dentate gyrus (mouse, neural) surfaces
\textsc{Gliogenesis}, \textsc{Neurogenesis}, and
\textsc{Central Nervous System Development}
(adj.\,$p$ down to $10^{-21}$);
hematopoiesis (mouse, blood) recovers
\textsc{Leukocyte Migration}, \textsc{Hemopoiesis}, and
\textsc{Granulocyte Activation};
Setty HSPC (human, immune) yields
\textsc{B Cell Activation} and \textsc{Lymphocyte Activation};
hESC (human, stem/hemostatic) captures
\textsc{Hemostasis}, \textsc{Platelet Activation}, and
\textsc{Wound Healing}
(adj.\,$p$ down to $10^{-23}$).
This tissue specificity, discovered \emph{without any supervision},
indicates that GAHIB learns disentangled, biologically meaningful
programmes that reflect underlying cellular biology rather than
generic transcriptional noise.

\begin{figure}[!htbp]
\centering
\includegraphics[width=\textwidth]{fig_downstream_analysis}
\caption{Unified downstream analysis of GAHIB.
Panels~(a)--(e) validate GAHIB's Lorentz-norm pseudotime from the
original gene expression on four datasets.
(a)~2D UMAP coloured by Leiden clusters.
(b)~Same embedding coloured by GAHIB pseudotime
(viridis: dark~$\to$~stem, light~$\to$~differentiated).
(c)~Scatter of GAHIB pseudotime vs. the top negatively-correlated
HVG (progenitor-like; expression decreases with $r_i$),
with the Pearson~$r$ annotated.
(d)~Scatter of GAHIB pseudotime vs. the top positively-correlated
HVG (lineage-commitment; expression increases with $r_i$).
(e)~Mean $|\text{Pearson}~r|$ of the top-20 HVGs per dataset,
summarising how strongly GAHIB's pseudotime tracks the dominant
gene programme.
Panels~(f)--(i) show tissue-specific GO~BP enrichment of the
ten latent dimensions for four independently trained GAHIB models,
demonstrating that latent dimensions specialise to
tissue-appropriate programmes without cross-context leakage.}
\label{fig:downstream}
\end{figure}


\subsection{Gene Attribution and Cross-Dataset Summary}

GAHIB recovers on average 31.5\% of DE marker genes within its
top attributed genes (\cref{tab:interp}), and identifies 24.7
significant GO enrichment terms per dataset among gene programs
aligned with individual latent dimensions.

\Cref{fig:interp_summary} provides an integrated overview combining
top gene attribution (panel~a) with stemness projections (panel~b).
GAT homophily is consistently high ($>0.75$, mean 0.812), confirming
that the attention mechanism learns biologically meaningful cell--cell
interactions.

\Cref{tab:interp} summarises key interpretation metrics per dataset.

\begin{figure}[!htbp]
\centering
\includegraphics[width=\textwidth]{interp_summary}
% Source: GAHIB_results/interpretation/figures/fig_themed_gene_stemness.pdf
% Generator: gahib/viz/interpretation.py::fig_themed_gene_stemness()
\caption{Gene attribution and stemness summary across all 53 datasets.
\textbf{(a)}~UMAP coloured by the top attributed gene (decoder
Jacobian $|\partial g/\partial z|$) per dataset; coherent spatial
gradients indicate that GAHIB captures biologically meaningful gene
programs.
\textbf{(b)}~Stemness score projected on UMAP
(red$=$high stemness, blue$=$low), showing the
stem$\to$differentiated gradient recovered by the hyperbolic latent
geometry.}
\label{fig:interp_summary}
\end{figure}

% --- Interpretation Summary Table ---
\begin{table}[!htbp]
\centering
\caption{Per-dataset interpretation summary (12 representative datasets shown;
means in text are computed over all 53 datasets).
Hom.\ = GAT attention homophily;
$\bar{r}_L$ = mean Lorentz norm;
$\rho_s$ = stemness--norm Spearman correlation
(\sigA/$p < 0.001$, except Pituitary: \sigC);
MO = marker overlap fraction;
GO = significant GO enrichment terms.}
\label{tab:interp}
\begin{tabular}{l c c c c c}
\toprule
Dataset & Hom. & $\bar{r}_L$ & $\rho_s$ & MO & GO \\
\midrule
TcellLiver   & 0.845 & 4.14 & 0.120\sigA & 0.426 & 46 \\
MCCTumor     & 0.861 & 1.17 & 0.501\sigA & 0.460 & 50 \\
bcEC         & 0.798 & 1.66 & $-$0.099\sigA & 0.358 & 33 \\
CA           & 0.808 & 1.89 & 0.505\sigA & 0.255 & 23 \\
Hepatobl.    & 0.869 & 1.81 & 0.275\sigA & 0.218 & 31 \\
Breast       & 0.851 & 1.84 & 0.615\sigA & 0.315 & 0 \\
Endo         & 0.878 & 1.66 & 0.386\sigA & 0.350 & 0 \\
Pituitary    & 0.789 & 2.31 & 0.043\sigC & 0.193 & 30 \\
Setty        & 0.860 & 1.46 & 0.179\sigA & 0.407 & 21 \\
hESC         & 0.758 & 1.73 & 0.454\sigA & 0.363 & 41 \\
Lung         & 0.827 & 1.37 & 0.274\sigA & 0.295 & 21 \\
Dentate      & 0.874 & 1.71 & 0.673\sigA & 0.402 & 0 \\
\midrule
\textbf{Mean (12)} & \textbf{0.835} & \textbf{1.90} & \textbf{0.327} & \textbf{0.337} & \textbf{24.7} \\
\bottomrule
\end{tabular}
\end{table}


% ============================================================================
\section{Robustness, Efficiency, and Sensitivity}\label{sec:robustness}
% ============================================================================

To address the concerns most commonly raised during peer review of
deep learning models in bioinformatics, we provide four additional
studies covering latent dimension selection, reproducibility across
random seeds, computational cost, and hyperparameter sensitivity.
All studies reuse the 53 scRNA-seq datasets and the same 20-metric
evaluation protocol.

\subsection{Latent Dimension Ablation}\label{sec:latdim}

We sweep the latent dimension $d\in\{3,5,10,20,50\}$ while holding the
IB dimension fixed at~2 and all other hyperparameters at their default
values.
\Cref{fig:latdim} reports mean performance across all 53 datasets.
Clustering quality peaks at $d{=}3$ (NMI\,$=$\,0.652, ASW\,$=$\,0.418)
and degrades monotonically as $d$ increases, with $d{=}50$ producing
markedly lower ASW (0.221) due to distance-concentration effects in
high-dimensional Euclidean space.
We retain $d{=}10$ as the default because it preserves sufficient
capacity for downstream tasks (trajectory inference, gene attribution)
while remaining within the stable region.

\begin{figure}[!htbp]
\centering
\includegraphics[width=0.95\textwidth]{latent_dim_ablation}
\caption{Latent dimension ablation across 53 datasets.
Bars show mean NMI, ARI, and ASW; error bars denote one standard
deviation.
The default $d{=}10$ (red outline) lies near the stable plateau
before ASW degrades at $d{=}50$.}
\label{fig:latdim}
\end{figure}

\subsection{Multi-Seed Reproducibility}\label{sec:seeds}

We retrain GAHIB with five random seeds
(\{42, 123, 456, 789, 2024\}) on each of the 53 datasets.
\Cref{fig:seeds} shows per-dataset mean$\pm$std NMI, ARI, and ASW.
Aggregate statistics are NMI\,$=$\,0.739$\pm$0.022,
ARI\,$=$\,0.546$\pm$0.041, and ASW\,$=$\,0.359$\pm$0.020,
confirming that GAHIB is highly reproducible: the within-dataset
standard deviation across seeds is typically less than 3\% of the mean.

\begin{figure}[!htbp]
\centering
\includegraphics[width=0.95\textwidth]{seed_robustness}
\caption{Multi-seed reproducibility analysis.
Each horizontal bar denotes one dataset (53 total);
error bars show the standard deviation across 5 random seeds.
Red dashed lines mark the overall mean.}
\label{fig:seeds}
\end{figure}

\subsection{Computational Cost and Scaling}\label{sec:cost}

\Cref{fig:cost}(a) compares training time and peak GPU memory for
three architectures on an NVIDIA RTX 4080 Laptop GPU.
GAHIB requires 22.0$\pm$3.2\,s per dataset and 0.11\,GB memory,
a 41\% overhead over Base VAE (15.6$\pm$2.4\,s, 0.055\,GB) and
10\% over VAE+IB+Hyp (20.0\,s, 0.059\,GB).
The additional cost comes primarily from subgraph sampling and
multi-head GAT message passing; GPU memory remains well below
production-scale limits.

\Cref{fig:cost}(b) examines scaling behaviour by training on
progressively larger subsamples (500, 1000, 2000, 3000 cells) of ten
datasets.
Training time is remarkably flat (22.0--22.6\,s across all cell counts)
because the bottleneck is the fixed number of subgraph training steps
(10 subgraphs $\times$ 512 nodes per epoch), not the dataset size.
This indicates that GAHIB's wall-clock cost does not grow linearly with
the number of cells, a useful property for atlas-scale applications.

\begin{figure}[!htbp]
\centering
\includegraphics[width=0.95\textwidth]{computational_cost}
\caption{Computational cost analysis.
\textbf{(a)}~Training time (bars, mean$\pm$std) and peak GPU memory
(annotations) for three methods across 53 datasets.
\textbf{(b)}~GAHIB training time as a function of the number of cells,
measured on 10 representative datasets at
$\{500, 1000, 2000, 3000\}$ cells.}
\label{fig:cost}
\end{figure}

\subsection{Hyperparameter Sensitivity}\label{sec:hpsens}

We sweep each of four key hyperparameters, holding the others at
default values: the KL weight
$\beta \in \{0.01, 0.05, 0.1, 0.5, 1.0\}$, the IB weight
$\lambda_\text{ib} \in \{0.1, 0.25, 0.5, 1.0, 2.0\}$, the Lorentz weight
$\lambda_\text{hyp} \in \{1.0, 2.5, 5.0, 10.0, 20.0\}$, and the graph
size $k \in \{5, 10, 15, 20, 30\}$.
\Cref{fig:hpsens} shows that GAHIB is robust to all four:
\begin{itemize}[leftmargin=*,itemsep=1pt]
\item $\beta$: NMI varies within a narrow range of
  $[0.723, 0.741]$ (total spread $\Delta{=}0.019$).
\item $\lambda_\text{ib}$: $[0.733, 0.745]$ ($\Delta{=}0.012$).
\item $\lambda_\text{hyp}$: $[0.711, 0.755]$ ($\Delta{=}0.044$).
\item $k$: $[0.718, 0.744]$ ($\Delta{=}0.027$).
\end{itemize}
The Lorentz weight shows the largest effect: smaller
$\lambda_\text{hyp}$ improves clustering (NMI\,$=$\,0.755 at
$\lambda_\text{hyp}{=}1$) but correspondingly reduces hierarchical
structure in the latent space, demonstrating the inherent trade-off
between cluster separation and hyperbolic geometry preservation.
The default $\lambda_\text{hyp}{=}5$ balances the two objectives.

\begin{figure}[!htbp]
\centering
\includegraphics[width=\textwidth]{hyperparam_sensitivity}
\caption{Hyperparameter sensitivity across the five key evaluation
metrics (NMI, ARI, ASW, DRE UMAP overall quality, LSE overall quality).
Each column varies a single hyperparameter while holding others fixed.
Red stars indicate the default value used in all other experiments.
GAHIB is robust across all tested ranges, with $\lambda_\text{hyp}$
showing the largest influence.}
\label{fig:hpsens}
\end{figure}


% ============================================================================
\section{Discussion}\label{sec:discussion}
% ============================================================================

\subsection{GAHIB's Main Practical Gain Is Balance}

The most important practical result is not that GAHIB wins every single
metric in isolation, but that it is the only method that remains strong
across clustering, geometry-preservation, and latent-structure analysis
at the same time.
The benchmark suite shows that competing methods each dominate only part
of the problem.
scVI comes close to GAHIB on NMI and ARI in the deep-learning benchmark,
scDHMap is competitive on geometry-aware metrics, and PCA or truncated
SVD remain strong on datasets with mostly linear structure.
In real single-cell workflows, however, users rarely train separate
embeddings for clustering, visualisation, pseudotime, and interpretation;
they want one representation that supports all of these downstream uses
without manual switching between models.

GAHIB is valuable precisely because it reduces that trade-off.
Compared with scVI, it preserves substantially stronger cluster margins
and latent geometry; compared with scDHMap, it improves separation while
keeping the hierarchy signal; compared with classical methods, it is more
reliable when tumour heterogeneity and developmental continua coexist in
the same dataset.
This suggests that GAHIB should be viewed less as a specialised
hyperbolic model and more as a balanced general-purpose embedding model
for scRNA-seq settings where both neighbourhood fidelity and branching
organisation matter.

\subsection{The GAT Encoder as the Critical Component}

Across all seven experimental studies, the GAT encoder consistently emerges
as the single most impactful architectural component.
The ablation (\cref{sec:ablation}) reveals that replacing the MLP encoder
with GAT (VAE+IB+Hyp $\to$ GAHIB) yields the largest performance gain on
every clustering metric (NMI +0.072, ARI +0.078), surpassing the effect
of adding either the information bottleneck or hyperbolic geometry individually.
The encoder comparison (\cref{sec:encoder}) further confirms a clear
and statistically significant hierarchy: GAT $>$ Transformer $>$ MLP.

GAHIB's superiority on \emph{cluster-separation} metrics---ASW (+0.149
vs.\ VAE+IB+Hyp), CAL, and DAV---follows directly from GAT-based
neighbourhood aggregation.
Mean attention homophily of 0.812 indicates that the encoder preferentially
propagates information across biologically similar neighbours, thereby
structuring the representation before variational compression.
This neighbourhood-aware encoding yields lower intra-cluster variance,
higher inter-cluster separation, and a more stable posterior geometry for
both NB reconstruction and IB compression.

GAHIB also leads on DRE-local metrics ($Q_\text{local}$, distance
correlation), reflecting that the neighbourhood structure of the
cell graph is faithfully preserved in the latent space.
Cells that are adjacent in the 15-NN graph retain similar latent positions,
so graph topology is encoded as continuous local proximity rather than as a
purely post-hoc clustering signal.

This result underscores a broader principle: explicitly encoding
cell-neighbourhood structure during representation learning provides a
stronger inductive bias than any form of posterior regularisation or
geometric loss applied post-hoc.
Accordingly, future single-cell VAEs should prioritise graph-aware encoders
before introducing increasingly elaborate latent-space regularisers.

\subsection{The IB--Hyperbolic Synergy}

A central architectural property of GAHIB is the {\em functional coupling}
between the information bottleneck and the hyperbolic geometry loss.
The IB layer compresses the $d$-dimensional latent code $\bz$ into a
2D manifold coordinate~$\bz'$; this compressed representation then
serves as the {\em anchor} for the Lorentz distance in
\cref{eq:hyp_loss}.
Without this anchor, the hyperbolic loss operates on an untrained
target and becomes degenerate.
The ablation confirms this directly:
VAE+Hyp (without IB) performs nearly identically to the Base VAE on
clustering (NMI 0.653 vs 0.651), whereas VAE+IB+Hyp recovers the full
DRE and LSE benefit (DRE 0.735, LSE 0.534).
IB enables the hyperbolic term, while the hyperbolic term in turn imposes
hierarchical organisation on the compressed manifold.
This reciprocal interaction provides a principled route for incorporating
geometric inductive bias into deep generative models.

\subsection{Hyperbolic Geometry Encodes Biological Hierarchy}

The interpretation study (\cref{sec:interpretation}) provides
quantitative evidence that GAHIB's Lorentz norm faithfully tracks
developmental hierarchy.
Stemness--norm Spearman correlations are positive and statistically
significant ($p < 0.001$) in the majority of datasets, with the strongest
associations observed in developmental systems (dentate gyrus:
$\rho_s = 0.673$; breast cancer: $0.615$; colorectal adenocarcinoma:
$0.505$; Merkel cell carcinoma: $0.501$).
The Poincar\'{e} disk projections (\cref{fig:interp_hyperbolic})
provide corresponding visual confirmation: stem and progenitor
populations cluster near the disk origin, while differentiated lineages
radiate toward the boundary.

The biological rationale for this behaviour follows from the geometry of
the Lorentz hyperboloid~\cite{ganea2018hyperbolic}.
The origin of hyperbolic space is the unique point of minimum norm; in tree
embeddings~\cite{nickel2017poincare}, the root of the hierarchy maps to
the origin while leaves map to high-norm positions.
GAHIB's geodesic loss encourages the high-dimensional latent code $\bz$
and its 2D projection $\bz'$ to occupy consistent positions on the
hyperboloid, effectively regularising the encoder to produce representations
whose \emph{radial distance from the origin encodes developmental
maturity}.
Stem and progenitor cells, which transcriptionally resemble the common
ancestor of many lineages, are naturally placed near the origin; their
descendants, which have progressively narrowed gene-expression programmes,
are pushed toward the boundary.
This mapping is learned entirely unsupervised---without pseudotime labels,
lineage trees, or any prior knowledge of developmental order.

These findings show that the Lorentz geometry loss does more than improve
aggregate metrics: it recovers a biologically coherent radial ordering of
cell states and provides a continuous \emph{stemness coordinate} as an
emergent property of the variational objective.

\subsection{Why Disentanglement Regularisation Hurts Single-Cell Clustering}

The disentanglement benchmark (\cref{sec:benchmark_disent}) reveals a
consistent and striking pattern: every disentanglement regulariser tested
(\(\beta\)-VAE, DIP-VAE, TC-VAE, InfoVAE) \emph{degrades} clustering quality
relative to the unregularised Base VAE, with TC-VAE producing the most severe
collapse (NMI 0.404, $-$0.243 vs.\ Base VAE).

Disentanglement objectives penalise statistical dependence among latent
dimensions, encouraging a posterior distribution whose marginals are as
independent as possible~\cite{higgins2017betavae,chen2018isolating}.
In natural images, this is beneficial because independent factors (shape,
colour, pose) genuinely exist and each dimension can specialise.
In scRNA-seq data, however, the dominant sources of variation are
co-regulated gene programmes---transcription factor networks, cell-cycle
machinery, lineage-specific regulons---which are fundamentally
\emph{correlated} by construction.
Forcing the encoder to represent these programmes in orthogonal dimensions
requires sacrificing reconstruction fidelity (higher ELBO cost), which
degrades the quality of the latent code for all downstream tasks.

InfoVAE~\cite{zhao2019infovae}, which relaxes the factorisation constraint,
is almost neutral (NMI 0.649 vs.\ Base 0.647), consistent with this
interpretation: minimal factorisation pressure preserves correlation
structure.
GAHIB's IB layer, by contrast, does not impose independence---it imposes a
\emph{geometric} constraint (compact 2D projection with small reconstruction
error), which selects for a low-dimensional manifold without disrupting the
correlation structure of the full 10D latent code.
This distinction---geometric compression rather than statistical
independence---explains why GAHIB simultaneously sustains strong clustering
(NMI 0.722) and high-fidelity dimensionality reduction (DRE 0.700).

\subsection{Biological Interpretability as a Design Principle}

GAHIB's architectural choices produce directly interpretable biological
signals.
The GAT encoder's attention weights encode cell-type homophily
(mean 0.812), indicating that the model assigns greater influence to
same-type neighbours without explicit supervision.
The Lorentz norm's correlation with stemness ($\rho_s = 0.327$ mean;
up to 0.673 in dentate gyrus) further demonstrates that hyperbolic
geometry captures the continuous nature of differentiation rather than
treating it as a discrete cluster assignment.
Together, these observations indicate that inductive biases aligned with
known biological structure can yield unsupervised biological alignment
without requiring cell-type labels during training.

\subsection{Latent Dimensions Encode Context-Specific Biology}

Gene Ontology Biological Process enrichment on four tissue contexts
(\cref{sec:go}) provides perhaps the most stringent biological test
of GAHIB's latent space: when trained independently on each dataset,
the model rediscovers \emph{the} biology of each tissue without any
cross-contamination.
The brain-derived dentate gyrus model returns \textsc{Gliogenesis},
\textsc{Neurogenesis}, and \textsc{CNS Development}
(adj.~$p{\leq}10^{-14}$);
the hematopoietic model returns
\textsc{Leukocyte Migration}, \textsc{Hemopoiesis}, and
\textsc{Granulocyte Activation};
and so on, each with the relevant gene markers appearing as expected
(Sox9, Olig2 for glia; Spi1, Elane for myeloid cells).
Critically, a latent dimension from the brain-trained model is never
enriched for blood biology and vice versa---%
a compositional property that disentanglement regularisers do not
achieve because they optimise statistical independence, not biological
meaningfulness.
This result confirms that the combined GAT$+$IB$+$Hyperbolic framework
yields latent factors that are simultaneously (i) statistically informative
for clustering, and (ii) biologically coherent to the trained tissue.

\subsection{Reproducibility, Cost, and Practical Deployability}

A deep-learning tool for single-cell analysis is only useful if it is
reproducible, robust, and efficient.
Multi-seed experiments (\cref{sec:seeds}) show that
GAHIB's clustering performance varies by less than 3\% across 5 random
seeds, addressing the common concern that deep-learning pipelines can
yield divergent results on identical data.
Hyperparameter sensitivity (\cref{sec:hpsens}) reveals that only
$\lambda_\text{hyp}$ has a meaningful effect on performance
(NMI spread $\Delta{=}0.044$), and even this effect is monotonic and
easily tuned; all four hyperparameter ranges include the default value
within the stable plateau.
Computational cost (\cref{sec:cost}) shows 22\,s per dataset on a
mid-range RTX 4080 Laptop GPU with 0.11\,GB peak memory---%
a modest 41\% overhead over Base VAE for a substantial performance
gain.
Scaling is notably flat with cell count because the training loop
samples a fixed number of subgraphs per epoch, a design choice that
should allow GAHIB to handle atlas-scale datasets (${>}100$K cells)
without retraining.

\subsection{Limitations and Future Directions}

\textbf{Computational cost.}
GAHIB requires constructing a $k$-NN graph and performing multi-head
GAT message passing, adding approximately 50\% wall-clock overhead
relative to an MLP-based VAE.
For interactive exploratory analysis on small cohorts, this is negligible;
for large production pipelines, graph construction can be pre-computed
once and cached.

\textbf{Scalability.}
Our evaluation caps at 3{,}000 cells per dataset to ensure fair
comparison across all methods.
Scaling to atlas-level data ($>$100K cells) would require efficient
graph-sampling strategies (e.g., GraphSAINT, ClusterGCN) and
mini-batch graph construction, which we leave to future work.
The subgraph sampling strategy we already employ (10 subgraphs of 512
nodes per epoch) provides a natural foundation for such extensions.

\textbf{Classical DR competitiveness on linear data.}
PCA achieves near-identical NMI/ARI to GAHIB on several datasets with
linearly separable structure (e.g., NMI 0.735 vs.\ 0.721), indicating
that GAHIB's advantage is largest for datasets with nonlinear
developmental or tumour heterogeneity.
Future work could explore hybrid approaches that adaptively select
between linear and graph-based encoders based on intrinsic
dimensionality estimates.

\textbf{Cell-type label quality.}
Our evaluation relies on Leiden clustering labels as ground truth.
While this provides a consistent unsupervised reference, it may not
perfectly capture biologically validated cell types.
Evaluating against manually curated annotations on benchmark datasets
would provide stronger biological validation.


% ============================================================================
\section{Conclusion}\label{sec:conclusion}
% ============================================================================

We presented GAHIB, a graph attention VAE with hyperbolic information
bottleneck for single-cell omics latent space learning.
The architecture unifies three synergistic inductive biases---GAT
encoding, information bottleneck compression, and Lorentz hyperbolic
geometry---into a coherent framework where each component enables the
others.

Through eleven experimental studies on 53 scRNA-seq datasets, we
demonstrated that:
\begin{enumerate}[leftmargin=*,itemsep=2pt]
\item Each component contributes to performance, with the GAT encoder
  providing the largest gain and the IB layer being architecturally
  necessary for hyperbolic geometry.
\item GAHIB achieves the strongest overall clustering among 7 deep
  learning and 5 classical baselines, with statistically significant
  gains in most pairwise comparisons.
\item GAHIB's latent space is biologically interpretable: hyperbolic radii
  track developmental hierarchy ($\rho_s = 0.327$ mean), decoder
  Jacobians recover 31.5\% of marker genes, and GAT attention exhibits
  81.2\% cell-type homophily.
\item GAHIB is robust, reproducible, and efficient: performance varies
  by less than 3\% across 5 random seeds, remains stable across
  hyperparameter sweeps (NMI spread $\leq 0.044$), and trains in
  22\,s/dataset with flat scaling across cell counts.
\end{enumerate}

These results establish that coupling graph-structured encoding with
hyperbolic geometry through an information bottleneck yields a principled
and empirically effective framework for biologically meaningful single-cell
latent space learning. More broadly, the results indicate that neighbourhood
structure and hierarchical geometry should be treated as complementary,
rather than competing, inductive biases in future single-cell generative
models.


% ============================================================================
\section*{Data and Code Availability}
% ============================================================================

\paragraph{Code.}
All source code, trained models, experiment scripts, and figure
generators are publicly available under an open-source licence at
\url{https://github.com/PeterPonyu/GAHIB}.
The repository includes the 11 experiment runners, the visualisation
pipeline, and reproducibility instructions.

\paragraph{Datasets.}
The 53 scRNA-seq datasets are all publicly available from the NCBI
Gene Expression Omnibus (GEO) under the following 42 unique
accession numbers:
GSE98638, GSE115571, GSE117988, GSE120446, GSE120505, GSE123813,
GSE123902, GSE124310, GSE130148, GSE132509, GSE138709, GSE142653,
GSE143423, GSE144024, GSE145929, GSE148215, GSE148218, GSE149655,
GSE155109, GSE163558, GSE165784, GSE165844, GSE167597, GSE168181,
GSE183904, GSE189070, GSE189357, GSE192857, GSE213740, GSE222002,
GSE222369, GSE225600, GSE225857, GSE226131, GSE226824, GSE228499,
GSE235787, GSE247719, GSE253355, GSE262288, GSE275119, GSE283205.
Together with the accessions listed above, the evaluation includes
five additional public benchmark datasets obtained from the
scDHMap benchmark repository: dentate gyrus, endoderm (endo),
hematopoiesis (hemato), Setty HSPC (setty), and lung, bringing the
total to 47 distinct data sources covering the 53 preprocessed
evaluation datasets (several GEO accessions contain multiple
independent subsets, e.g.\ GSE123813 provides both basal and
squamous cell carcinoma cohorts).
Preprocessing scripts that reproduce the exact evaluation splits are
included in the code repository.


% ============================================================================
\section*{Author Contributions}
% ============================================================================

\textbf{Zeyu Fu:} conceptualization, methodology, software, formal
analysis, investigation, data curation, visualization, writing---original
draft, writing---review and editing.


% ============================================================================
\section*{Funding}
% ============================================================================

This research received no specific grant from any funding agency in the
public, commercial, or not-for-profit sectors.


% ============================================================================
\section*{Declaration of Competing Interest}
% ============================================================================

The author declares no competing financial or personal interests that
could have appeared to influence the work reported in this paper.


% ============================================================================
\section*{Ethics Statement}
% ============================================================================

This study uses only de-identified, publicly available scRNA-seq data
obtained from the NCBI Gene Expression Omnibus (GEO).
All data were released by the original authors under open-access
licences and ethical approvals were obtained by the original data
producers as indicated in their respective publications.
No new human or animal data were collected in this study.


% ============================================================================
\section*{Declaration of AI Usage}
% ============================================================================

The author used AI-based writing-assistance tools (large language models)
during manuscript preparation for grammar checking, wording refinement,
and literature search.
All scientific contributions, analyses, experimental designs,
interpretations, and conclusions are the work of the author.
The author reviewed and takes full responsibility for all content.


% ============================================================================
\bibliographystyle{elsarticle-num}
\bibliography{references}


\end{document}
